% GENERATED FILE. Edit canonical lessons, then run npm run generate:books.
% canonical-source-hash: fnv1a64:ae32004d23ba5e9f
% canonical-lessons: HI-C472-phurtila, HI-C472-sust, HI-C472-nikamma, HI-C472-karmath, HI-C472-mustaid

\chapter{Lazy or Brisk}
\label{ch:hi-lazy-or-brisk}
\hlchaptermodality{472}

\begin{chapteropening}
\textbf{By the end of this chapter:} \emph{I can say whether someone gets things done.}

Complete the last lesson of chapter 472: I can say whether someone gets things done.
\end{chapteropening}

\section[\texorpdfstring{\hi{फुर्तीला} (phurtīlā)}{फुर्तीला (phurtīlā)}]{\hi{फुर्तीला} (phurtīlā) — nimble; brisk}

\label{lesson:HI-C472-phurtila}



\begin{quote}
Before the new one: say the Hindi for etiquette, then the Hindi for to tease.
\end{quote}

\subsection*{You'll want to know: \hi{फुर्तीला}}
\textbf{\hi{फुर्तीला}} — \emph{phurtīlā} — \textquotedblleft{}nimble; brisk\textquotedblright{}.

\begin{quote}
\hi{फुर्तीला} \hi{खिलाड़ी} \hi{गेंद} \hi{को} \hi{जल्दी} \hi{पकड़} \hi{लेता} \hi{है।}
\end{quote}

\begin{quote}
\emph{phurtīlā khilāṛī gẽd ko jaldī pakaṛ letā hai.}
\end{quote}

\begin{quote}
A nimble player catches the ball quickly.
\end{quote}

\begin{cousinweb}
\textbf{\hi{फुर्तीला}} (\emph{phurtīlā}) is \emph{phurtī}, \textquotedblleft{}briskness\textquotedblright{}, with the ending \emph{-īlā}, \textquotedblleft{}full of\textquotedblright{}; CDIAL lists Hindi \emph{phurtī}, \textquotedblleft{}activity\textquotedblright{}, under Sanskrit \emph{sphuráti}, \textquotedblleft{}darts, quivers\textquotedblright{} (CDIAL 13849).
\end{cousinweb}

\begin{grammarlens}[title={What you've built}]
One more word for describing what you lived and felt.
\end{grammarlens}

\subsection*{Your turn}
\begin{itemize}
\raggedright
  \item \emph{Say it:} \emph{phurtīlā}
  \item \emph{Say it:} \emph{phurtīlā}, once more
  \item \emph{Say it:} \emph{phurtīlā khilāṛī gẽd ko jaldī pakaṛ letā hai.} — A nimble player catches the ball quickly.
  \item \emph{Recall:} say the Hindi for etiquette, then the Hindi for to tease, then say \emph{phurtīlā} again
\end{itemize}

\begin{tcolorbox}[breakable,colback=teal!4,colframe=teal!35!black,title={Before you move on}]
What does \hi{फुर्तीला} mean? (\textquotedblleft{}Nimble; brisk\textquotedblright{}.) Say it once more.
\end{tcolorbox}
\section[\texorpdfstring{\hi{सुस्त} (sust)}{सुस्त (sust)}]{\hi{सुस्त} (sust) — sluggish; slack}

\label{lesson:HI-C472-sust}



\begin{quote}
Before the new one: say the Hindi for to tease, then the Hindi for nimble.
\end{quote}

\subsection*{You'll want to know: \hi{सुस्त}}
\textbf{\hi{सुस्त}} — \emph{sust} — \textquotedblleft{}sluggish; slack\textquotedblright{}.

\begin{quote}
\hi{गर्मी} \hi{में} \hi{सब} \hi{लोग} \hi{सुस्त} \hi{हो} \hi{जाते} \hi{हैं।}
\end{quote}

\begin{quote}
\emph{garmī meṁ sab log sust ho jāte haiṁ.}
\end{quote}

\begin{quote}
In the heat everyone becomes sluggish.
\end{quote}

\begin{cousinweb}
\textbf{\hi{सुस्त}} (\emph{sust}) is usually traced to Persian \emph{sust}, \textquotedblleft{}weak, slack\textquotedblright{}.
\end{cousinweb}

\begin{grammarlens}[title={What you've built}]
One more word for describing what you lived and felt.
\end{grammarlens}

\subsection*{Your turn}
\begin{itemize}
\raggedright
  \item \emph{Say it:} \emph{sust}
  \item \emph{Say it:} \emph{sust}, once more
  \item \emph{Say it:} \emph{garmī meṁ sab log sust ho jāte haiṁ.} — In the heat everyone becomes sluggish.
  \item \emph{Recall:} say the Hindi for to tease, then the Hindi for nimble, then say \emph{sust} again
\end{itemize}

\begin{tcolorbox}[breakable,colback=teal!4,colframe=teal!35!black,title={Before you move on}]
What does \hi{सुस्त} mean? (\textquotedblleft{}Sluggish; slack\textquotedblright{}.) Say it once more.
\end{tcolorbox}
\section[\texorpdfstring{\hi{निकम्मा} (nikammā)}{निकम्मा (nikammā)}]{\hi{निकम्मा} (nikammā) — good-for-nothing; idle}

\label{lesson:HI-C472-nikamma}



\begin{quote}
Before the new one: say the Hindi for nimble, then the Hindi for sluggish.
\end{quote}

\subsection*{You'll want to know: \hi{निकम्मा}}
\textbf{\hi{निकम्मा}} — \emph{nikammā} — \textquotedblleft{}good-for-nothing; idle\textquotedblright{}.

\begin{quote}
\hi{वह} \hi{लड़का} \hi{निकम्मा} \hi{है}, \hi{दिन} \hi{भर} \hi{कुछ} \hi{नहीं} \hi{करता।}
\end{quote}

\begin{quote}
\emph{vah laṛkā nikammā hai, din bhar kuch nahīṃ kartā.}
\end{quote}

\begin{quote}
That boy is good for nothing; he does nothing all day.
\end{quote}

\begin{cousinweb}
\textbf{\hi{निकम्मा}} (\emph{nikammā}) is a tadbhava: CDIAL lists Hindi \emph{nikammā}, \textquotedblleft{}useless, unemployed\textquotedblright{}, under Sanskrit \emph{niṣkarman} (CDIAL 7475), from \emph{niḥ}, \textquotedblleft{}without\textquotedblright{}, and \emph{karman}, \textquotedblleft{}work\textquotedblright{}.
\end{cousinweb}

\begin{grammarlens}[title={What you've built}]
One more word for describing what you lived and felt.
\end{grammarlens}

\subsection*{Your turn}
\begin{itemize}
\raggedright
  \item \emph{Say it:} \emph{nikammā}
  \item \emph{Say it:} \emph{nikammā}, once more
  \item \emph{Say it:} \emph{vah laṛkā nikammā hai, din bhar kuch nahīṃ kartā.} — That boy is good for nothing; he does nothing all day.
  \item \emph{Recall:} say the Hindi for nimble, then the Hindi for sluggish, then say \emph{nikammā} again
\end{itemize}

\begin{tcolorbox}[breakable,colback=teal!4,colframe=teal!35!black,title={Before you move on}]
What does \hi{निकम्मा} mean? (\textquotedblleft{}Good-for-nothing; idle\textquotedblright{}.) Say it once more.
\end{tcolorbox}
\section[\texorpdfstring{\hi{कर्मठ} (karmaṭh)}{कर्मठ (karmaṭh)}]{\hi{कर्मठ} (karmaṭh) — dedicated; diligent in work}

\label{lesson:HI-C472-karmath}



\begin{quote}
Before the new one: say the Hindi for sluggish, then the Hindi for good-for-nothing.
\end{quote}

\subsection*{You'll want to know: \hi{कर्मठ}}
\textbf{\hi{कर्मठ}} — \emph{karmaṭh} — \textquotedblleft{}dedicated; diligent in work\textquotedblright{}.

\begin{quote}
\hi{मीरा} \hi{की} \hi{माँ} \hi{बहुत} \hi{कर्मठ} \hi{हैं}, \hi{सुबह} \hi{से} \hi{रात} \hi{तक} \hi{काम} \hi{करती} \hi{हैं।}
\end{quote}

\begin{quote}
\emph{Mīrā kī māṁ bahut karmaṭh haiṁ, subah se rāt tak kām kartī haiṁ.}
\end{quote}

\begin{quote}
Mira's mother is very dedicated; she works from morning till night.
\end{quote}

\begin{cousinweb}
\textbf{\hi{कर्मठ}} (\emph{karmaṭh}) is a tatsama, borrowed from Sanskrit \emph{karmaṭha}, \textquotedblleft{}capable of work; working diligently\textquotedblright{} (Monier-Williams), from \emph{karman}, \textquotedblleft{}work\textquotedblright{}.
\end{cousinweb}

\begin{grammarlens}[title={What you've built}]
One more word for describing what you lived and felt.
\end{grammarlens}

\subsection*{Your turn}
\begin{itemize}
\raggedright
  \item \emph{Say it:} \emph{karmaṭh}
  \item \emph{Say it:} \emph{karmaṭh}, once more
  \item \emph{Say it:} \emph{Mīrā kī māṁ bahut karmaṭh haiṁ, subah se rāt tak kām kartī haiṁ.} — Mira's mother is very dedicated; she works from morning till night.
  \item \emph{Recall:} say the Hindi for sluggish, then the Hindi for good-for-nothing, then say \emph{karmaṭh} again
\end{itemize}

\begin{tcolorbox}[breakable,colback=teal!4,colframe=teal!35!black,title={Before you move on}]
What does \hi{कर्मठ} mean? (\textquotedblleft{}Dedicated; diligent in work\textquotedblright{}.) Say it once more.
\end{tcolorbox}
\section[\texorpdfstring{\hi{मुस्तैद} (mustaid)}{मुस्तैद (mustaid)}]{\hi{मुस्तैद} (mustaid) — alert; ready for action}

\label{lesson:HI-C472-mustaid}



\begin{quote}
Before the new one: say the Hindi for good-for-nothing, then the Hindi for dedicated.
\end{quote}

\subsection*{You'll want to know: \hi{मुस्तैद}}
\textbf{\hi{मुस्तैद}} — \emph{mustaid} — \textquotedblleft{}alert; ready for action\textquotedblright{}.

\begin{quote}
\hi{चुनाव} \hi{के} \hi{दिन} \hi{पुलिस} \hi{मुस्तैद} \hi{थी।}
\end{quote}

\begin{quote}
\emph{cunāv ke din pulis mustaid thī.}
\end{quote}

\begin{quote}
On election day the police were on the alert.
\end{quote}

\begin{cousinweb}
\textbf{\hi{मुस्तैद}} (\emph{mustaid}) is usually traced to Arabic \emph{mustaʿidd}, \textquotedblleft{}ready, prepared\textquotedblright{}, which reached Hindi through Persian.
\end{cousinweb}

\begin{grammarlens}[title={What you've built}]
One more word for describing what you lived and felt.
\end{grammarlens}

\subsection*{Your turn}
\begin{itemize}
\raggedright
  \item \emph{Say it:} \emph{mustaid}
  \item \emph{Say it:} \emph{mustaid}, once more
  \item \emph{Say it:} \emph{cunāv ke din pulis mustaid thī.} — On election day the police were on the alert.
  \item \emph{Recall:} say the Hindi for good-for-nothing, then the Hindi for dedicated, then say \emph{mustaid} again
\end{itemize}

\begin{tcolorbox}[breakable,colback=teal!4,colframe=teal!35!black,title={Before you move on}]
What does \hi{मुस्तैद} mean? (\textquotedblleft{}Alert; ready for action\textquotedblright{}.) Say it once more.
\end{tcolorbox}
